% ===== 底本 第 1 章《引言》；片段 tr01a；源块 B012--B023 =====

\section{引言}

% [源 B012]
在本次报告中，我想从更偏数学的角度，对 SYZ 镜像对称纲领作一个简要概述。回忆 Kontsevich [Kon94] 的\emph{同调镜像对称}（HMS）纲领：它断言存在成对的（未必光滑的）簇 $(X,\omega,J)$ 与 $(\check{X},\check{\omega},\check{J})$，使得
\begin{align*}
  &D^{\pi}\mathcal{Fuk}(X,\omega)\simeq D^b\operatorname{Coh}(\check{X},\check{J})\\
  &\qquad\text{且}\quad D^b\operatorname{Coh}(X,J)\simeq D^{\pi}\mathcal{Fuk}(\check{X},\check{\omega}),
\end{align*}
其中 $\mathcal{Fuk}(X,\omega)$ 是辛流形 $(X,\omega)$ 的某个版本的“Fukaya 范畴”，$D^b\operatorname{Coh}(X,J)$ 是复簇 $(X,J)$ 的有界凝聚层导出范畴，只要它们有良好定义。在代数一侧（我们称之为\emph{B-模型}），有 Bondal 与 Orlov [BO01] 的一个著名定理：我们可以从其凝聚层导出范畴 $D^b\operatorname{Coh}(X)$ 重构出光滑射影簇 $X$。

% [Tutor Note: B012 首现 “Homological Mirror Symmetry”，按术语表定译「同调镜像对称（HMS）」；式子中的 $D^{\pi}$ 是分裂闭包（Karoubian 完备化）后的导出 Fukaya 范畴，源文以 $\pi$ 作上标，照抄不改。]

% [源 B013]
\begin{theorem*}[Bondal--Orlov]
设 $X$ 是光滑射影簇，$Y$ 是代数簇，且典范层 $K_X$ 或反典范层 $-K_X$ 丰沛，则对任何导出范畴的等价
\[
  D^b\operatorname{Coh}(X)\xrightarrow[\cong]{f}D^b\operatorname{Coh}(Y),
\]
都存在概形同构 $\varphi\colon X\xrightarrow{\cong}Y$，使得 $D^b(\varphi_{\ast})=f$。
\end{theorem*}

% [Tutor Note: B013 的前提只把 $Y$ 称作代数簇，但结论给出概形同构 $\varphi\colon X\xrightarrow{\cong}Y$，故 $Y$ 的光滑性与射影性自动成立；此处照译原文，并非漏译条件。]

% [源 B014]
\begin{proof}[证明梗概]
由于这一证明思路对 SYZ 图景很重要，我们简要说明该定理证明的想法。首先，设 $X$ 与 $Y$ 都是射影簇，其坐标代数分别为 $B_X=\bigoplus_{i=0}^{\infty}H^i(X,\mathcal{O}_X(i))$ 与 $B_Y=\bigoplus_{i=0}^{\infty}H^i(Y,\mathcal{O}_Y(i))$，则任何分次同构 $B_X\cong B_Y$ 都会诱导相应簇的同构 $X\cong Y$。既然已知 $X$ 是射影簇，且其典范除子或反典范除子 $\pm K_X$ 丰沛，就只需从导出范畴的同构 $D^b\operatorname{Coh}(X)\xrightarrow{f}\mathcal{D}$ 中找出 $\mathcal{O}_X(1)$（等价地，$\pm K_X$）的像，并证明它在 $Y$ 中可由一个丰沛除子代表。为此，我们需要一个非交换意义下的“丰沛除子”概念，以便在抽象导出范畴 $\mathcal{D}$ 中把相应的对象辨认出来。在带有\emph{Serre 函子}的抽象三角范畴 $\mathcal{D}$ 中---所谓 Serre 函子，即一个自同构 $S_{\mathcal{D}}\colon\mathcal{D}\to\mathcal{D}$，使得对任意 $A,B\in\mathcal{D}$ 都有 $\operatorname{Hom}_{\mathcal{D}}(A,B)\cong\operatorname{Hom}_{\mathcal{D}}(B,SA)^{\ast}$---我们可以如下定义\emph{点对象}：称对象 $P$ 是余维数为 $s$ 的点对象，如果
\begin{enumerate}[label=\roman*.,leftmargin=1.9em,itemsep=1pt,topsep=2pt]
% [源 B015]
  \item $S_{\mathcal{D}}(P)\simeq P[s]$；
% [源 B016]
  \item $\operatorname{Hom}^{\leq 0}(P,P)=0$；
% [源 B017]
  \item $\operatorname{Hom}^0(P,P)=k(P)$。
\end{enumerate}

% [Tutor Note: B014 中 Serre 函子的对偶写作 $\operatorname{Hom}_{\mathcal{D}}(B,SA)^{\ast}$，$SA$ 是 $S_{\mathcal{D}}A$ 的简写；核对时可取 $\mathcal{D}=D^b\operatorname{Coh}(X)$、$S=-\otimes\omega_X[\dim X]$，此时 $S(\mathcal{O}_p)\simeq\mathcal{O}_p[n]$ 正合 B015 的 $P[s]$（$s=n=\dim X$）。]

% [源 B018]
其中 $k(P)$ 是 $k$ 的一个扩域。显然，$X$ 上的摩天层自然地是 $D^b\operatorname{Coh}(X)\simeq\mathcal{D}$ 中的点对象；反之，$D^b\operatorname{Coh}(X)$ 中的任何点对象都是 $X$ 上摩天层的平移。由于这是一个纯范畴的概念，$\mathcal{D}$ 中点对象的同构类与 $D^b\operatorname{Coh}(X)$ 中点对象的同构类一一对应。上述论证只给出了簇 $X$ 与 $Y$ 的闭点集之间的等同，我们还需要等同相应的拓扑以及概形结构。拓扑（以及最终的同构）的关键要素是\emph{线束}的范畴式概念，或用代数几何的语言，\emph{可逆层}。称 $\mathcal{D}$ 中的对象为\emph{可逆的}，如果对任何点对象 $P\in\mathcal{D}$，都能找到 $s\in\mathbb{Z}$，使得
% [Tutor Note: B018 的 $k(P)$ 是点对象的自同态代数 $\operatorname{Hom}^0(P,P)$，原文说它是 $k$ 的扩域：对 $X$ 上的摩天层 $\mathcal{O}_p$，$k(P)$ 就是剩余域 $\kappa(p)$；B017 与 B019/B020 中的 $k(P)$ 必须理解为同一个量。]
\begin{enumerate}[label=\roman*.,leftmargin=1.9em,itemsep=1pt,topsep=2pt]
% [源 B019]
  \item $\operatorname{Hom}^s(L,P)=k(P)$；
% [源 B020]
  \item $\operatorname{Hom}^i(L,P)=0$，对 $i\neq s$。
\end{enumerate}

% [源 B021]
还可以证明，若 $\mathcal{D}$ 来自某个簇的导出范畴，则这类对象至多相差一个平移就同构于某个线束。现在固定 $X$ 上的线束 $L_0$，考虑满足 $\operatorname{Hom}(L_0,P)=k(P)$ 的点对象 $P$。这些对象将是 $Y$ 中未经平移的摩天层；而对任意一对可逆层 $L_1,L_2$ 及态射 $\alpha\colon L_1\to L_2$，诱导态射
\[
  \Hom(L_2,P)\xrightarrow{\alpha_P^{\ast}}\Hom(L_1,P)
\]
定义出一族点对象 $U_{\alpha}=\{P\vert\alpha_P^{\ast}\neq 0\}$。由于 $X$ 是射影簇，当 $\alpha$ 取遍所有满足 $\Hom^0(L_i,P)=k(P)$ 的可逆对象 $L_1,L_2$ 时，这些 $U_{\alpha}$ 构成 $X$ 的 Zariski 拓扑的一组基。正如对任何代数簇一样，这样的开集总构成相应 Zariski 拓扑的一组基 [ND20]，这实际上表明 $X$ 与 $Y$ 作为拓扑空间是同胚的。剩下的任务是确定典范层/反典范层，并证明簇 $Y$ 也是 Fano 型/一般型的。对固定的可逆层 $L_0$，可考虑线束 $L_i=S^iL_0[-ni]$，则这些 $L_i$ 也是 $X$ 中的线束（在 $Y$ 中亦然），而代数 $\bigoplus_{i=0}^{\infty}\Hom(L_0,L_i)$ 就是 $X$ 上典范层的自同态代数，在 $Y$ 上同样如此。现在，由于 $\omega_X$ 或是丰沛的或是反丰沛的，对所有 $i,j$、取 $\alpha\in\Hom(L_i,L_j)$ 得到的开集 $\{U_{\alpha}\}$ 就足以构成 Zariski 拓扑的一组基；而这是一个当且仅当的条件，因此相应的典范层 $\omega_Y$ 也应当是丰沛的或反丰沛的。至此，开头一段中的论证就成立了，于是得到同构 $X\cong Y$。

% [Tutor Note: B021 末句原文作 “should also be ample and anti-ample”，与同段 “$\omega_X$ is either ample or anti-ample” 相矛盾（丰沛与反丰沛互斥），语义应为「或」；译文按「或」处理，已列入回执。另同段 $K_X$ 与 $\omega_X$ 混用，指同一个典范层。]
\end{proof}

% [源 B022]
基于重构定理，只要代数簇是 Fano 型或一般型，我们就能从它的凝聚层导出范畴粗略地“重构”出这个代数簇。我们暂且忽略这一条件，把它当作启发式的原则。现在我们从辛几何得到了范畴的一个来源：对每个辛流形 $(X,\omega)$ 我们都可以设想一个 $A_{\infty}$-范畴 $\mathcal{Fuk}(X,\omega)$；若镜像存在，它应当成为某个代数簇 $\check{X}$ 的凝聚层导出范畴的 $A_{\infty}$-增强，特别地，我们应当能从范畴 $\mathcal{Fuk}(X,\omega)$ 重构出 $X$。假定 $\check{X}$ 是分离的，则由 [ND20] 我们只需确定点对象与可逆对象。

% [Tutor Note: B022 末句 “we should be able to re-construct $X$ from the category $\mathcal{Fuk}(X,\omega)$” 里的 $X$ 按上下文应为镜像 $\check{X}$：下一句才以「假定 $\check{X}$ 分离」为前提。此处照译存疑，已列入回执。]

% [Tutor Note: B022 源文作 "resconstruction theorem"（多一个 s），应为 reconstruction；译文按正字处理，不改数学。]

% [源 B023]
设 $p\in\check{X}$ 是代数簇 $(\check{X},\check{J})$ 的光滑闭点，$\mathcal{O}_p$ 是闭点 $p$ 的结构层。则可以直接计算得到：

% [Tutor Note: B023 句末 “we can directly compute that” 后面直接接 B024 的引理陈述（中间无独立公式），故译文用冒号收束；引理及其证明从 tr01b 开始，属同一章内跨片段的句子衔接，不影响环境完整性。]
